\chapter{Validation of the Coupled System}
\label{chap:coupled-validation}

In this chapter, we study the coupled motion of an incompressible fluid and an
inextensible Euler--Bernoulli fibre. The fluid is advanced by the adaptive
projection solver discussed in Section~\ref{sec:centered-projection-solver},
while the fibre is advanced by the dynamic ADDM beam solver described in
Chapter~\ref{chap:addm}.

\section{Coupled Formulation}

Let \(\Omega\) denote the fluid domain and let \(\vec
r:[0,L]\times[0,T]\to\mathbb{R}^3\) denote the fibre centerline. We write
\(\bm\Lambda(s,t)\) for the hydrodynamic force density exerted by the fluid on
the fibre. The equal and opposite force is applied to the fluid through the
regularized immersed-boundary kernel \(\delta_h\). A formal continuous
statement of the coupled system is
\begin{align}
    \rho_f\left(
    \pdv{\vec u}{t}
    + \vec u\cdot\nabla\vec u
    \right)
     & =
    -\nabla p
    + \mu_f\nabla^2\vec u
    + \rho_f\vec g
    \notag                               \\
     & \quad -
    \int_0^L
    \bm\Lambda(s,t)
    \delta_h(\vec x-\vec r(s,t))
    \odif{s},
    \label{eq:coupled-fluid-momentum}    \\
    \nabla\cdot\vec u
     & = 0,
    \label{eq:coupled-incompressibility} \\
    \rho_b\pdv[2]{\vec r}{t}
    + EI\pdv[4]{\vec r}{s}
    - \pdv*{\left(T\pdv{\vec r}{s}\right)}{s}
     & =
    \vec f_b + \bm\Lambda,
    \label{eq:coupled-beam-momentum}     \\
    \norm{\pdv{\vec r}{s}}^2
     & = 1,
    \label{eq:coupled-inextensibility}   \\
    \pdv{\vec r}{t}(s,t)
     & =
    \int_\Omega
    \vec u(\vec x,t)
    \delta_h(\vec x-\vec r(s,t))
    \odif{\vec x}.
    \label{eq:coupled-noslip}
\end{align}
Here \(T\) is the fibre tension enforcing inextensibility, and \(\vec f_b\)
denotes any non-hydrodynamic body force applied directly to the beam. The last
equation is the immersed no-slip condition: the Lagrangian fibre velocity is
the Eulerian fluid velocity interpolated to the fibre.

The numerical method is a non-boundary-fitted, partitioned immersed-boundary
coupling. The fluid variables live on the adaptive Eulerian grid, while the
fibre is represented by Lagrangian marker points on the beam centerline. At a
time step, the immersed-boundary forcing procedure computes a marker force
density associated with the residual slip between the fluid and fibre. This
force is spread to the fluid with the opposite sign and is also passed as the
hydrodynamic load to the ADDM beam solver. The beam solve then advances the
inextensible centerline and returns updated marker positions and velocities to
the fluid solver.

This coupling keeps the fluid and beam solvers modular and avoids remeshing
under large fibre motion. Its constraints are enforced by different pieces of
the algorithm: pressure projection enforces fluid incompressibility, the ADDM
beam iteration enforces fibre inextensibility, and the immersed-boundary force
reduces the no-slip residual at the Lagrangian markers. The method is therefore
not a monolithic solve of
\eqref{eq:coupled-fluid-momentum}--\eqref{eq:coupled-noslip}; the coupling
accuracy is limited by the regularized transfer kernel, the immersed-boundary
forcing iteration, and the time-step splitting.

Because the immersed-boundary force is computed before the final pressure
projection, the final projected velocity does not satisfy the marker no-slip
constraint exactly. The coupling should therefore be understood as a
partitioned, weakly coupled time-splitting scheme, with no-slip enforced up to
the accuracy of the forcing iteration, transfer kernel, and projection split.

\section{Validation}

\subsection{Single Flapping Fibre in a Uniform Flow}

We compare our results to a case first presented in
\textcite{huang_simulation_2007} and in subsequent studies
\cite{wang_combined_2008,xin_efficient_2023,goza_strongly-coupled_2016,lee_discrete-forcing_2015,wang_strongly_2015}.

\begin{figure}[t]
    \centering
    \def\svgwidth{0.5\textwidth}
    \import{figure/vector}{huang-validation-domain.pdf_tex}
    \caption{Physical domain setup for the single flapping fibre in a uniform flow.}
    \label{fig:flapping-fibre-validation}
\end{figure}

The reference problem consists of a single inextensible filament with one end
fixed and the other end free, placed in a uniform incoming flow. Huang \emph{et
al.} nondimensionalize the equations using the filament length \(L_r\), the
incoming speed \(U_\infty\), and the fluid density \(\rho_f\). In these
variables, the fluid equations take the form
\begin{equation}
    \pdv{\vec u}{t}
    + \vec u \cdot \nabla \vec u
    =
    -\nabla p
    + \frac{1}{\Reynolds}\nabla^2 \vec u
    + \vec f,
    \qquad
    \nabla \cdot \vec u = 0,
\end{equation}
where
\begin{equation}
    \Reynolds = \frac{\rho_f U_\infty L_r}{\mu_f}.
\end{equation}
With the sign convention used in this thesis, the corresponding dimensionless
beam equation is
\begin{equation}
    \pdv[2]{\vec r}{t}
    =
    \pdv{}{s}\left( T\pdv{\vec r}{s} \right)
    - \pdv[2]{}{s}\left( K_B\pdv[2]{\vec r}{s} \right)
    + {\rm Fr}\,\hat{\vec g}
    + \bm\Lambda,
    \qquad
    \pdv{\vec r}{s}\cdot\pdv{\vec r}{s} = 1 .
\end{equation}
Here \(K_B\) is the dimensionless bending rigidity, \(T\) is the tension
enforcing inextensibility, and \(\bm\Lambda\) denotes the hydrodynamic force
density exerted on the beam. Following Huang \emph{et al.}, the gravitational
parameter is \({\rm Fr}=gL_r/U_\infty^2\), which is the inverse-square form of
the more common Froude number. The density parameter reported in the reference
case is \(q=\rho_b/(\rho_f L_r)\), the dimensionless beam line density.

For the single-fibre validation, the initially straight filament is inclined by
\(\theta_0=0.1\pi\). The fixed end is simply supported and the free end has
zero tension, bending moment, and shear force. Huang \emph{et al.} used the
domain \([-2,6]\times[-4,4]\) in coordinates scaled by \(L_r\). Our
implementation uses the same nondimensional parameter set in a larger
computational box, with the inlet placed \(2L_r\) upstream of the fixed end.
The comparison is therefore based on the same physical benchmark, while the
no-slip constraint is imposed by the immersed-boundary coupling described above
rather than by Huang's feedback-force parameters.

\begin{table}[H]
    \centering
    \begin{tabular}{ll}
        Quantity                        & Value           \\ \hline
        Reynolds number, \(\Reynolds\)  & \(200\)         \\
        Density ratio, \(q\)            & \(1.5\)         \\
        Bending rigidity, \(K_B\)       & \(10^{-3}\)     \\
        Gravity parameter, \({\rm Fr}\) & \(0.5\)         \\
        Filament length, \(L/L_r\)      & \(1\)           \\
        Initial angle, \(\theta_0\)     & \(0.1\pi\)      \\
        Beam resolution                 & \(64\) elements \\
    \end{tabular}
    \caption{Dimensionless parameters for the single flapping fibre validation
        case of \textcite{huang_simulation_2007}.}
    \label{tab:huang-single-fibre-setup}
\end{table}

\begin{figure}[H]
    \centering
    \begin{tabular}{@{}cccc@{}}
        \includegraphics[width=0.23\textwidth]{figure/raster/single-fibre-Re-200-KB-10e-3_early_t9p2.png}  &
        \includegraphics[width=0.23\textwidth]{figure/raster/single-fibre-Re-200-KB-10e-3_early_t10.png}   &
        \includegraphics[width=0.23\textwidth]{figure/raster/single-fibre-Re-200-KB-10e-3_early_t10p8.png} &
        \includegraphics[width=0.23\textwidth]{figure/raster/single-fibre-Re-200-KB-10e-3_early_t11p6.png}   \\
        \(t=9.2\)                                                                                          &
        \(t=10.0\)                                                                                         &
        \(t=10.8\)                                                                                         &
        \(t=11.6\)
    \end{tabular}
    \caption{Snapshots from the \(\Reynolds=200\) single-fibre validation case with \(K_B = 10^{-3}\).
        Black curves show vorticity contours and red markers show the Lagrangian
        fibre points.}
    \label{fig:single-fibre-re-200-kb-10e-3-snapshots}
\end{figure}

\begin{table}[H]
    \centering
    \begin{tabular}{c|cc}
                                              & Amplitude    & Frequency (St) \\ \hline
        \textcite{huang_simulation_2007}      & \(\pm 0.35\) & \(0.30\)       \\
        \textcite{wang_strongly_2015}         & \(\pm 0.35\) & \(0.31\)       \\
        \textcite{lee_discrete-forcing_2015}  & \(\pm 0.38\) & \(0.31\)       \\
        \textcite{goza_strongly-coupled_2016} & \(\pm 0.38\) & \(0.32\)       \\
        Present                               & \(\pm 0.36\) & \(0.31\)       \\
    \end{tabular}
    \caption{Amplitude and frequency associated with the transverse displacement
        of the beam's trailing edge; obtained for \(\Reynolds = 200\), \(q = 1.5\),
        \(K_B = 10^{-3}\) and \({\rm Fr} = 0.5\).}
    \label{tab:huang-re-200-kb-10e-3-results}
\end{table}

\begin{figure}[t]
    \centering
    \import{figure/vector}{huang-tip-Re-200-KB-10e-4.tex}
    \caption{Time history of the free end of the filament at with nondimensional parameters from Table~\ref{tab:huang-single-fibre-setup} and \(\Delta t = 0.0003\). In dashed lines, results from \textcite{huang_simulation_2007}.}
    \label{fig:huang-tip-re-200-KB-10e-4}
\end{figure}

\subsection{Single Settling Fibre in a Quiescent Fluid}

We
